%% summary.tex
\section{总结与展望}

\begin{frame}{工作总结与展望}
    \scriptsize
    \begin{columns}
        \begin{block}{总结}
            \vspace{0.5em}
            改进版AutoPT+总体成功率从原版的24.67\%提升至63.67\%（+39\%），同时单次测试成本降低69\%、耗时缩短42\%，其中对弱模型gpt-3.5-turbo提升尤为显著。
        \end{block}   
        \begin{alertblock}{现有不足}
            \begin{enumerate}
                \item \textbf{提前放弃（33\%）}：重试建议没有区分失败类型
                \item \textbf{安全审查拦截（22\%）}：模型拒绝生成攻击性命令
                \item \textbf{覆盖范围局限}：仅针对 Web 应用层漏洞
            \end{enumerate}
        \end{alertblock}
        \begin{block}{未来工作}
            \begin{itemize}
                \item 精细化重试策略（按失败类型）
                \item 引入本地开源模型（如 DeepSeek）
                \item 扩展 PSM 覆盖完整渗透测试阶段
            \end{itemize}
        \end{block}
    \end{columns}
\end{frame}
